\subsection{\texorpdfstring{\textcolor{red}{Steps 1--3: }}{Steps 1--3: }Preprocessing and exact low-degree certification}
\label{app:upper-proof-screen}

We first establish the residual Markov and tail certificates, then prove the
fresh-color success bound and the low-degree screen contract.
\textcolor{red}{These certify the vocabulary banks from
\AlgoBlockNamedRef{preprocess}, the fresh colors from
\AlgoBlockNamedRef{draft}, and the rows returned by
\AlgoBlockNamedRef{screen}.}

The next lemma derives the residual Markov property from
\cref{main:eq:forest-factorization}.
Recall from \eqref{app:eq:history} that
$\ResidualPositions=\PositionSet\setminus G$ and
$\ResidualForest=\TargetForest[\ResidualPositions]$.
\textcolor{red}{Its row identity is used in \AlgoBlockNamedRef{probes},
and its component independence justifies \AlgoBlockNamedRef{centroid}.}

\begin{lemma}[Residual Markov property]
\label{app:lem:residual-markov}
Fix a history $\History=(G,x_G)$.  Conditional on $X_G=x_G$, the variables in
distinct connected components of $\ResidualForest$ are independent.  Moreover,
fix $j\in\ResidualPositions$ and the values of all residual neighbors of $j$.
The exact conditional row of $X_j$ is unchanged if any residual nonneighbor is
revealed, masked, changed, or marginalized.
\end{lemma}

\begin{proof}

\emph{1. Factor the conditional law over residual components.}
Recall the forest factorization \eqref{app:eq:S0}:
$\TargetLaw(x)=\prod_{u\in\TargetRoots}\PosMarginal{u}(x_u)
 \prod_{j\notin\TargetRoots}\EndpointKernel{\ParentOf{j}}{j}
 (x_j\mid x_{\ParentOf{j}})$.
Here $\TargetRoots$ and $\ParentOf{\cdot}$ are the original reference roots
and parent map.  Substitute the fixed committed values $x_G$.
For each connected component $T$ of $\ResidualForest$, identified with
its vertex set, define
\[
\begin{aligned}
 g_T(x_T):={}&
 \prod_{u\in\TargetRoots\cap T}\PosMarginal{u}(x_u)
 &&\text{roots in }T\\
 &{}\times
 \prod_{\substack{j\in T\setminus\TargetRoots\\\ParentOf{j}\in T}}
 \EndpointKernel{\ParentOf{j}}{j}(x_j\mid x_{\ParentOf{j}})
 &&\text{edges within }T\\
 &{}\times
 \prod_{\substack{j\in T\setminus\TargetRoots\\\ParentOf{j}\in G}}
 \EndpointKernel{\ParentOf{j}}{j}(x_j\mid x_{\ParentOf{j}})
 &&\text{edges }G\to T\\
 &{}\times
 \prod_{\substack{j\in G\setminus\TargetRoots\\\ParentOf{j}\in T}}
 \EndpointKernel{\ParentOf{j}}{j}(x_j\mid x_{\ParentOf{j}})
 &&\text{edges }T\to G.
\end{aligned}
\]
The fixed $x_G$ dependence is suppressed in $g_T$.
In particular, the last product retains factors for committed children of
vertices in $T$: their parent values still vary with $x_T$.
Every edge with a residual endpoint is either internal to one such $T$ or
crosses between $T$ and $G$; no edge joins distinct residual components.
Thus, with products over $T$ ranging over these connected components,
\[
 \TargetLaw(x_G,x_{\ResidualPositions})
 =
 \left[
 \prod_{u\in\TargetRoots\cap G}\PosMarginal{u}(x_u)
 \prod_{\substack{j\in G\setminus\TargetRoots\\\ParentOf{j}\in G}}
 \EndpointKernel{\ParentOf{j}}{j}(x_j\mid x_{\ParentOf{j}})
 \right]\prod_T g_T(x_T).
\]
The bracket is positive and depends only on $x_G$.
It therefore cancels when dividing this joint probability by
$\TargetLaw(X_G=x_G)
 =\sum_{z_{\ResidualPositions}}\TargetLaw(x_G,z_{\ResidualPositions})$.
Since the connected components partition $\ResidualPositions$,
finite distributivity gives
\[
 \sum_{x_{\ResidualPositions}}\prod_T g_T(x_T)
 =\prod_T\left(\sum_{x_T}g_T(x_T)\right).
\]
Consequently
\[
\begin{aligned}
 \TargetLaw(X_{\ResidualPositions}=x_{\ResidualPositions}\mid X_G=x_G)
 &=\frac{\prod_T g_T(x_T)}
         {\sum_{z_{\ResidualPositions}}\prod_T g_T(z_T)}\\
 &=\prod_T\frac{g_T(x_T)}{\sum_{z_T}g_T(z_T)}.
\end{aligned}
\]
All sums range over the vocabulary assignments on their indicated vertex
sets; their denominators are positive by strict positivity of the factors.
This proves conditional independence of the residual components.


\emph{2. Cancel all factors not incident to the readout.}
For the second claim, root the original connected component at $j$,
using (C0).
Every original neighbor $v$ of $j$ is observed: either $v\in G$, or its
value is fixed in the lemma.  Call this value $x_v$.  The factors involving
$X_j=a$ are exactly
\[
 \PosMarginal{j}(a)
 \prod_{v\in\mathcal N_{\TargetForest}(j)}
       \EndpointKernel{j}{v}(x_v\mid a).
\]
All other factors are independent of $a$ once these neighbor values are
fixed.  Revealing, changing, or summing out nonneighbors only changes a
multiplicative factor independent of $a$.  It cancels in the normalized row
\[
 \frac{\PosMarginal{j}(a)
       \prod_{v\in\mathcal N_{\TargetForest}(j)}
          \EndpointKernel{j}{v}(x_v\mid a)}
      {\sum_{b\in\Vocab}\PosMarginal{j}(b)
       \prod_{v\in\mathcal N_{\TargetForest}(j)}
          \EndpointKernel{j}{v}(x_v\mid b)}.
\]
Thus the history and fixed neighbor values determine the row.
\end{proof}

The next lemma certifies the omitted tokens for screening and, under (RF\textcolor{red}{; \cref{app:eq:RF}})
with the standard floor, bounds the number of source columns.

\paragraph{Intuition.}
The local vocabulary bank $\TokenBank{i}$ from \eqref{app:eq:local-bank}
includes every token whose true marginal mass exceeds $\TailThreshold$.  Conversely, under the standard floor, every included
token has true mass at least $\TailThreshold/2$.
The rank envelope then bounds how many such tokens exist:
\[
 \frac{\TailThreshold}{2}
 \le\FrequencyConstant(\BankSize{i}^{-\RankTailExponent}+\VocabSize^{-1}),
 \quad
 \FrequencyConstant\VocabSize^{-1}\le\frac{\TailThreshold}{4}
 \quad\Longrightarrow\quad
 \BankSize{i}\le(4\FrequencyConstant/\TailThreshold)^{1/\RankTailExponent}.
\]
This is an upper bound on bank size; (RF\textcolor{red}{; \cref{app:eq:RF}}) does not require the envelope to
be attained.  The complementary tokens can all be represented by one
certified tail token in the response test.

\begin{lemma}[Preprocessing certificates\textcolor{red}{; \AlgoBlockRef{preprocess}}]
\label{app:lem:preprocessing-certificates}
Suppose preprocessing in \cref{app:def:preprocessing} is feasible, and let
$\TailThreshold$ be the threshold it selects.  Every local vocabulary bank
$\TokenBank{i}$ has a nonempty complement in $\Vocab$, so each
$\TailRepresentative{i}$ is well-defined.  For every
$i\in\PositionSet$ and $a\in\Vocab$,
\begin{equation}
 \bigl|\PosMarginal{i}(a)-\NoisyMarginal{i}(a)\bigr|
 \le\RowTVError,
 \qquad
 a\notin\TokenBank{i}
 \Longrightarrow a\in\TailSet{i}{\TailThreshold}.
 \label{app:eq:proof-preprocessing-certificates}
\end{equation}
Under (RF\textcolor{red}{; \cref{app:eq:RF}}), the standard floor \eqref{app:eq:rate-threshold-floor}
additionally gives
\begin{equation}
 \BankSize{i}
 \le\left(\frac{4\FrequencyConstant}{\TailThreshold}\right)^{1/\RankTailExponent}
 \quad\text{for every }i.
 \label{app:eq:finite-bank-size-proof}
\end{equation}
\end{lemma}

\begin{proof}

\emph{1. Establish a nonempty vocabulary-bank complement.}
In \AlgoBlockNamedRef{preprocess}, preprocessing fixes
$\TailThreshold=\max\FeasibleThresholdGrid$ once from the
public parameters, before the all-mask submission; this threshold is reused
by every screen.
Since $\FrequencyConstant\ge1$, the standard floor satisfies
$\ThresholdFloor=4\max\{\RowTVError,\FrequencyConstant/\VocabSize\}
\ge2(\RowTVError+1/\VocabSize)$.
Together with $\TailThreshold\ge\ThresholdFloor$, this gives
\begin{equation}
 \TailThreshold-\RowTVError
 \ge\RowTVError+\frac2{\VocabSize}
 >\frac1{\VocabSize}
 \ge\min_{a\in\Vocab}\NoisyMarginal{i}(a).
 \label{app:eq:baseline-existence}
\end{equation}
Thus $\TokenBank{i}\ne\Vocab$, proving tail-representative existence.

\emph{2. Certify every omitted token using the row-TV bound.}
At the all-mask state, $\ExactRow{\AllMaskState}{i}=\PosMarginal{i}$ and
$\OracleRow{\AllMaskState}{i}=\NoisyMarginal{i}$.
Thus \eqref{app:eq:oracle-tv-radius} and the singleton-event bound for TV give,
for every $a\in\Vocab$,
\[
 \bigl|\PosMarginal{i}(a)-\NoisyMarginal{i}(a)\bigr|
 \le \sup_{A\subseteq\Vocab}
       \bigl|\PosMarginal{i}(A)-\NoisyMarginal{i}(A)\bigr|
 =\dTV(\PosMarginal{i},\NoisyMarginal{i})
 \le\RowTVError.
\]
This bounds both signs of the difference, whether or not $a$ belongs to
$\TokenBank{i}$.  For a token outside the vocabulary bank $\TokenBank{i}$,
definition \eqref{app:eq:local-bank} then gives
\begin{equation}
 a\notin\TokenBank{i}
 \quad\Longrightarrow\quad
 \PosMarginal{i}(a)
 \le\NoisyMarginal{i}(a)+\RowTVError<\TailThreshold.
 \label{app:eq:bank-complement-tail}
\end{equation}
This proves \eqref{app:eq:proof-preprocessing-certificates}, including the
tail certificate for $\TailRepresentative{i}$.


\emph{3. Apply (RF\textcolor{red}{; \cref{app:eq:RF}}) only for the bank-size estimate.}
For the vocabulary-bank size $\BankSize{i}=|\TokenBank{i}|$,
the standard floor gives
$\TailThreshold\ge\ThresholdFloor=\StandardThresholdFloor{\FrequencyConstant}$,
hence both $\TailThreshold\ge4\RowTVError$ and
$\TailThreshold\ge4\FrequencyConstant/\VocabSize$.
Every $a\in\TokenBank{i}$ therefore satisfies
\begin{equation}
 \PosMarginal{i}(a)
 \ge\NoisyMarginal{i}(a)-\RowTVError
 \ge\TailThreshold-2\RowTVError
 \ge\frac{\TailThreshold}{2}.
 \label{app:eq:bank-member-mass}
\end{equation}
If
$\BankSize{i}\ge1$, at least $\BankSize{i}$ tokens have exact mass at least
$\TailThreshold/2$; hence the exact rank-$\BankSize{i}$ mass has the same
lower bound.  Assumption (RF\textcolor{red}{; \cref{app:eq:RF}}) and
$\TailThreshold\ge4\FrequencyConstant/\VocabSize$ yield
\[
 \frac{\TailThreshold}{2}
 \le\FrequencyConstant
 \left(\BankSize{i}^{-\RankTailExponent}+\VocabSize^{-1}\right)
 \le\FrequencyConstant\BankSize{i}^{-\RankTailExponent}
    +\frac{\TailThreshold}{4}.
\]
Rearranging proves \eqref{app:eq:finite-bank-size-proof}.  The inequality is
trivial when $\BankSize{i}=0$.
\end{proof}

\paragraph{Screen-success event (analysis only).}
\textcolor{red}{The colors drawn in \AlgoBlockNamedRef{draft} determine
the following event for \AlgoBlockNamedRef{screen}.}
Fix a realized history $\History=(G,x_G)$.  For distinct
$i,j\in\ResidualPositions$ with
$d_{\ResidualForest}(j)\le\DegreeCutoff$, call $\tau$
\emph{$(i,j)$-separating} when
\begin{equation}
\begin{aligned}
 \ColorHash{\tau}{v}&\ne\ColorHash{\tau}{j}
   &&\text{for every }v\in\mathcal N_{\ResidualForest}(j),\\
 \ColorHash{\tau}{v}&\ne\ColorHash{\tau}{i}
   &&\text{for every }
   v\in\mathcal N_{\ResidualForest}(j)\setminus\{i\}.
\end{aligned}
\label{app:eq:pair-separating-coloring}
\end{equation}
The screen family is \emph{successful at $\History$}, denoted
$\ScreenSuccess{\History}$, if
\begin{equation}
 \left|
 \left\{\tau\in\ColoringIndexSet:
       \tau\text{ is $(i,j)$-separating}\right\}
 \right|>\frac{\ColoringCount}{2}
 \label{app:eq:screen-success-event}
\end{equation}
for every such ordered pair $(i,j)$.  This event is defined using the hidden
residual forest only for analysis; the implemented screen neither observes
nor tests it.  For every separating $\tau$, the two properties used by the
screen proof are
\begin{equation}
 \mathcal N_{\ResidualForest}(j)
 \cap\ColorClass{\tau}{\ColorHash{\tau}{j}}=\emptyset,
 \qquad
 \mathcal N_{\ResidualForest}(j)
 \cap\ColorClass{\tau}{\ColorHash{\tau}{i}}\subseteq\{i\}.
 \label{app:eq:pairwise-local-isolation}
\end{equation}
\begin{example}[Checking which colors separate a pair]
\label{ex:algorithm-separation-20260918}
Use the coloring drawn in \AlgoBlockRef{draft} from
\cref{ex:algorithm-submissions-20260918} and the forest in
\cref{ex:algorithm-peel-centroids-20260918}.  For the tested pair
$(i,j)=(1,10)$, the readout has neighbors
$\mathcal N_{\ResidualForest}(10)=\{1,11\}$.  The two tests in
\eqref{app:eq:pairwise-local-isolation} are therefore
\[
 \{1,11\}\cap\ColorClass{\tau}{2}=\emptyset,\qquad
 \{1,11\}\cap\ColorClass{\tau}{1}=\{1\}\subseteq\{i\}.
\]
Both hold: $h_\tau(1)=1$, $h_\tau(10)=2$, and $h_\tau(11)=3$.
If instead $h_\tau(11)=2$, then neighbor $11$ shares the readout color
and the first intersection contains $11$.  If instead $h_\tau(11)=1$,
then an additional neighbor shares the source color and the second
intersection is $\{1,11\}\not\subseteq\{1\}$.
This checks one pair and one coloring.  It does not assert the
simultaneous strict-majority event \eqref{app:eq:screen-success-event},
which requires the full family of $\ColoringCount$ colorings.
\end{example}


\paragraph{Call index, draft, and information before fresh colors.}
Fix the target, its forest, and the frozen oracle throughout.
The index $s\in[\ScreenCallCap]$ counts screen calls
(\AlgoBlockRef{screen}), not individual colorings or commit rounds.
Define the \emph{reach event}
\[
 R_s:=\{\text{Algorithm~\ref{app:alg:complete-sampler} reaches its $s$th screen call}\}.
\]
Thus $R_s$ is an event, not the round count $\CommitRounds$.
Recall that $\DraftFiller=(\DraftToken{i})_{i\in\ResidualPositions}
\in\Vocab^{\ResidualPositions}$ is the background token vector fixed by
$\DraftGenerator$ before fresh colors are drawn
(Definition~\ref{app:def:draft-generator}).  Its call-indexed version is
\[
 f^{(s)}:=
 \begin{cases}
 \DraftFiller\text{ fixed for call }s\text{ in \AlgoBlockRef{draft}},
     &\text{on }R_s,\\
 \dagger,&\text{on }R_s^{\mathsf c}.
 \end{cases}
\]
The symbol $\dagger$ means \emph{no $s$th call}: it is a formal marker,
not a vocabulary token, $\MASK$, or a draft vector.
Different executions can stop after different numbers of screens.
This marker defines $f^{(s)}$ on every execution, so the proof can sum
over the fixed range $s=1,\ldots,\ScreenCallCap$ without adding any call.

On $R_s$, let $\mathcal T_{s-}$ record the past colors, replies, commits,
and decision-rule state just after $f^{(s)}$ is fixed and before the
current color block is drawn.  On $R_s^{\mathsf c}$, use the terminal
transcript instead.  Define
\[
 \mathcal F_{s-1}:=\sigma(\mathcal T_{s-},f^{(s)}),
 \qquad
 \mathcal F:=\mathcal F_{s-1}\quad\text{at the call under consideration}.
\]
On a reached call, this is precisely the information available between
draft generation and fresh color generation: the history and completed
draft are fixed under this conditioning, but the current colors are not.
It excludes current or future colors and the full seed $W$; fixing $W$
would fix those colors.  No union over possible histories or drafts is needed.

\paragraph{Intuition.}
A low-degree readout creates at most $2\DegreeCutoff$ forbidden color
equalities.  Thus
$\RandomColorCount=8(\DegreeCutoff+1)\simeq\DegreeCutoff$
makes one coloring good with probability greater than $3/4$.
Repeating independently $\RandomColoringCount$ times makes a wrong
majority exponentially unlikely:
\[
 \Pr(\hbox{wrong majority for one pair}\mid\mathcal F)
 \le e^{-\RandomColoringCount/8},\qquad
 \ScreenCallCap N^2 e^{-\RandomColoringCount/8}\le\ScreenFailureBudget.
\]
A union bound over pairs yields \cref{main:eq:conditional-screen-success}.
The chosen $\RandomColoringCount$ is the ceiling of the logarithm
obtained by solving the second inequality.  Conditioning is on the past before the fresh colors.

\begin{lemma}[Fresh randomized screens succeed adaptively\textcolor{red}{; \AlgoBlockRef{draft}}]
\label{app:lem:random-screen-success}
Suppose Algorithm~\ref{app:alg:complete-sampler} reaches a screen call with
realized past $\mathcal F$, including the completed generation of its fixed
draft but preceding the current coloring block.  Conditional on
$\mathcal F$, the fresh family in
\eqref{app:eq:random-color-parameters} satisfies
\begin{equation}
 \Pr\bigl(\ScreenSuccess{\History}^{\mathsf c}\mid\mathcal F\bigr)
 \le \frac{\ScreenFailureBudget}{\ScreenCallCap}.
 \label{app:eq:conditional-screen-failure}
\end{equation}
Consequently, over all decision-rule randomness and commit draws,
\begin{equation}
 \Pr\{\text{some reached screen is unsuccessful}\}
 \le\ScreenFailureBudget.
 \label{app:eq:adaptive-screen-failure}
\end{equation}
\end{lemma}

\begin{proof}

\emph{1. Separate one ordered pair with one fresh coloring.}
Conditional on $\mathcal F$, the history, residual forest, and draft are fixed, while
the current colors drawn in \AlgoBlockRef{draft} are mutually independent
and uniform.  Fix an ordered pair
$(i,j)$ with $d_{\ResidualForest}(j)\le\DegreeCutoff$.  The two lines of
\eqref{app:eq:pair-separating-coloring} contain at most
$2d_{\ResidualForest}(j)\le2\DegreeCutoff$ forbidden color equalities.  Each
has probability $1/\RandomColorCount$, so one coloring is nonseparating with
conditional probability at most
\begin{equation}
 \frac{2\DegreeCutoff}{\RandomColorCount}
 =\frac{\DegreeCutoff}{4(\DegreeCutoff+1)}<\frac14.
 \label{app:eq:one-random-coloring-bad}
\end{equation}

\emph{2. Amplify to a strict majority for that pair.}
Let $X_\tau$ indicate that coloring $\tau$ separates this pair.  Conditional
on $\mathcal F$, the $X_\tau$ are independent and have mean at least $3/4$.
Hoeffding's inequality therefore gives
\begin{equation}
 \Pr\left\{\sum_{\tau=1}^{\RandomColoringCount}X_\tau
       \le\frac{\RandomColoringCount}{2}\,\middle|\,\mathcal F\right\}
 \le \exp\{-2\RandomColoringCount(3/4-1/2)^2\}
 =\exp\{-\RandomColoringCount/8\}.
 \label{app:eq:random-majority-hoeffding}
\end{equation}

\emph{3. Take a union bound over ordered pairs at the current history.}
There are fewer than $N^2$ relevant ordered pairs.  A union bound and
\eqref{app:eq:random-color-parameters} give
\[
 N^2e^{-\RandomColoringCount/8}
 \le\ScreenFailureBudget/\ScreenCallCap,
\]
which proves \eqref{app:eq:conditional-screen-failure}.


\emph{4. Sum over adaptively reached calls by conditional expectation.}
Recall that $R_s$ means the algorithm reaches its $s$th screen call,
and $f^{(s)}$ is that call's draft; $\dagger$ marks an unreached call.
On $R_s$, let $H^{(s)}$ be the committed history at that call, and define
\hypertarget{app-reached-screen-failure}{}
\[
 E_s:=\{R_s\text{ occurs and }\ScreenSuccess{H^{(s)}}\text{ fails}\}.
\]
Thus $E_s$ means \emph{call $s$ is reached and fails}; it is false on
$R_s^{\mathsf c}$, with no history or screen evaluated there.
The draft rule terminates before the fresh colors
(Definition~\ref{app:def:draft-generator}), so reaching the call is
decided before those colors.  In the notation above,
$R_s=\{f^{(s)}\ne\dagger\}\in
\mathcal F_{s-1}=\sigma(\mathcal T_{s-},f^{(s)})$.
On $R_s$ apply \eqref{app:eq:conditional-screen-failure};
on $R_s^{\mathsf c}$ the conditional failure probability is zero.  Hence
\[
 \Pr(E_s\mid\mathcal F_{s-1})
 \le
 \begin{cases}
 \ScreenFailureBudget/\ScreenCallCap,&\text{on }R_s,\\
 0,&\text{on }R_s^{\mathsf c}
 \end{cases}
 =\mathbf 1_{R_s}\frac{\ScreenFailureBudget}{\ScreenCallCap}.
\]
Taking expectations gives
\[
 \Pr(E_s)
 =\mathbb E\!\left[
   \Pr(E_s\mid\mathcal F_{s-1})
  \right]
 \le\mathbb E\!\left[
   \mathbf 1_{R_s}\frac{\ScreenFailureBudget}{\ScreenCallCap}
  \right]
 =\Pr(R_s)\frac{\ScreenFailureBudget}{\ScreenCallCap}
 \le\frac{\ScreenFailureBudget}{\ScreenCallCap}.
\]
The failed-screen event is the union of these reached-call failures, so
\[
 \Pr\!\left(\bigcup_{s=1}^{\ScreenCallCap}E_s\right)
 \le\sum_{s=1}^{\ScreenCallCap}\Pr(E_s)
 \le\ScreenCallCap\,
      \frac{\ScreenFailureBudget}{\ScreenCallCap}
 =\ScreenFailureBudget.
\]
This proves \eqref{app:eq:adaptive-screen-failure}.  No independence between
different screen-success events is asserted or needed.  The same separating coloring
isolates every packed column for a fixed ordered pair, so neither the event
nor this union bound requires an additional union over bank values, columns,
or possible drafts.  Draft generation is allowed to depend on the past
oracle transcript; independence of that draft from oracle error is not used.
\end{proof}

\paragraph{Intuition\textcolor{red}{~(\AlgoBlockRef{chunks}, \AlgoBlockRef{probes})}.}
\textcolor{red}{By \cref{app:lem:random-screen-success}, with high probability
every reached screen has a strict separating majority for each pair with
readout degree at most $\DegreeCutoff$.
Fix one such coloring with different source and readout colors.
The readout chunk contains no neighbor of $j$, and the probed source class
contains no neighbor of $j$ other than possibly $i$
(\cref{app:eq:pairwise-local-isolation}).
Thus every neighbor is observed, and only $i$ can change among them.
By \cref{app:lem:residual-markov}, the exact row equals the row obtained by
varying only $i$ at the common draft boundary in
\cref{app:eq:fixed-draft-boundary-row}; this is the boundary at which
(RT; \cref{app:eq:RT}) and (UEN; \cref{app:eq:UEN}) apply.}

\begin{lemma}[Exact rows at a fixed draft boundary\textcolor{red}{; \AlgoBlockRef{chunks}, \AlgoBlockRef{probes}}]
\label{app:lem:fixed-draft-rows}
Suppose preprocessing is feasible and $\MaxBankSize\ge1$.
Fix $\History=(G,x_G)$, $\DraftFiller\in\Vocab^{\ResidualPositions}$, and
distinct residual positions $i,j$.  Define the complete boundary $z$ on
$\PositionSet\setminus\{i,j\}$ by $z_k=x_k$ for $k\in G$ and
$z_k=\DraftToken{k}$ otherwise.  For an $(i,j)$-separating coloring with
different source and readout colors, every $\kappa\in\ScreenColumnSet$
satisfies
\begin{equation}
 \ExactRow{
   \ProbeState{\kappa}{\tau}{\ColorHash{\tau}{i}}
              {\ReadoutChunkOf{\tau}{j}}}{j}
 =\BoundaryRow{i}{j}{\PackedColumn{\kappa}{i}}{z}.
 \label{app:eq:fixed-draft-boundary-row}
\end{equation}
In particular, the exact row family contains all source values in the
local vocabulary bank $\TokenBank{i}$ and its tail representative
$\TailRepresentative{i}$, at this same boundary.
Padding may additionally contribute the row for $\DraftToken{i}$.
\end{lemma}

\begin{proof}
By \AlgoBlockNamedRef{chunks}, the readout chunk is contained in the
readout color class.
Thus the two isolation clauses in \eqref{app:eq:pairwise-local-isolation}
give, respectively,
\[
 \mathcal N_{\ResidualForest}(j)\cap\ReadoutChunkOf{\tau}{j}
 =\varnothing,
 \qquad
 \mathcal N_{\ResidualForest}(j)\cap
       \ColorClass{\tau}{\ColorHash{\tau}{i}}
 \subseteq\{i\}.
\]
All residual neighbors of $j$ are therefore observed.  Since the source
and readout colors differ, $i$ is also outside the masked readout chunk.
In the state submitted by \AlgoBlockNamedRef{probes}, the relevant
observed coordinates are exactly
\[
 \left(\ProbeState{\kappa}{\tau}{\ColorHash{\tau}{i}}
                    {\ReadoutChunkOf{\tau}{j}}\right)_v
 =
 \begin{cases}
 x_v,&v\in G,\\
 \PackedColumn{\kappa}{i},&v=i,\\
 \DraftToken{v},&v\in\mathcal N_{\ResidualForest}(j)\setminus\{i\}.
 \end{cases}
\]
These are the same values as in the complete boundary with source value
$\PackedColumn{\kappa}{i}$ and remaining values $z$.  Conditional on
$X_G=x_G$, Lemma~\ref{app:lem:residual-markov} says that the row for $j$
depends only on its residual neighbors.  Completing or changing all other
coordinates to $z$ therefore leaves it unchanged, proving the identity.
The bank columns and column $\TailColumn$ give the asserted subfamily by
\eqref{app:eq:packed-column}.  The boundary is independent of the coloring
because one draft is fixed across the entire screen.
\end{proof}

\textcolor{red}{\Cref{app:lem:random-screen-success} bounds the probability
that any reached screen fails.  Under the screen-success event,
\cref{app:lem:low-degree-screen} formalizes \eqref{main:eq:exact-low-degree-row}
and identifies the output of
\AlgoBlockNamedRef{aggregate} with the true neighbor set for every readout
of residual degree at most $\DegreeCutoff$.}
The target obeys the forest and RT--UEN\textcolor{red}{~%
(\cref{app:eq:RT,app:eq:UEN})} assumptions
(Assumptions~\ref{app:ass:forest-structure}
and~\ref{app:ass:response-regularity}), and oracle rows satisfy the TV bound
\eqref{app:eq:oracle-tv-radius}.
Recall the general resolution
$\ScreenResolution=4\RowTVError+\TailTolerance$ from
\eqref{app:eq:diameter-screen-resolution}; the main text uses
$\RowTVError=\TargetTolerance/(2\sqrt N)$.
The lemma's screening-margin hypothesis $\ScreenResolution<\EdgeSignal$
means that the target signal exceeds the row-noise and omitted-response allowances.
This is a sufficient screening condition, not a consequence of RT--UEN\textcolor{red}{~%
(\cref{app:eq:RT,app:eq:UEN})} alone.
Preprocessing feasibility is specified in
Definition~\ref{app:def:preprocessing}, and
$\ScreenSuccess{\History}$ is defined in
\eqref{app:eq:screen-success-event}.

\begin{lemma}[Low-degree screen contract\textcolor{red}{; \AlgoBlockRef{aggregate}}]
\label{app:lem:low-degree-screen}
% Please retain the former label as a compatibility alias to this same lemma.
\label{main:lem:low-degree-screen}
Suppose preprocessing is feasible and
$\ScreenResolution<\EdgeSignal$ holds. At any history $\History$,
with any draft fixed throughout the screen, $\ScreenSuccess{\History}$
implies
\begin{equation*}
 \ScreenRow{j}=\mathcal N_{\ResidualForest}(j)
 \quad\text{for every }j\in\ResidualPositions
 \text{ with }d_{\ResidualForest}(j)\le\DegreeCutoff.
\end{equation*}
\end{lemma}
High-degree rows have only the construction's size cap; no subset or
correctness guarantee is asserted for them.

\paragraph{Intuition for the screen contract.}
At a separating coloring, all exact rows for a nonedge coincide, so its
observed diameter is at most $2\RowTVError$.  For an edge, restricting source
values to the vocabulary bank plus tail representative,
$\TokenBank{i}\cup\{\TailRepresentative{i}\}$, loses at most
$\TailTolerance$ of the full response
diameter, and row noise loses at most another $2\RowTVError$.  Thus the
edge's observed diameter is at least
\[
 \EdgeSignal-\TailTolerance-2\RowTVError>2\RowTVError.
\]
The strict inequality is exactly
$4\RowTVError+\TailTolerance<\EdgeSignal$.  The proof below verifies the
bank/tail comparison, all empty-bank cases, and the strict-majority step.

\begin{proof}[Proof of \AppLowDegreeScreenRef]

\emph{1. Fix a successful screen and the hidden exact rows.}
We analyze the votes and selected rows of \AlgoBlockNamedRef{aggregate}.
If $\MaxBankSize=0$, use the empty-bank argument in the final paragraph
below.  Otherwise assume $\MaxBankSize\ge1$.
Fix a history $\History$ and a screen draft $\DraftFiller$ for which
$\ScreenSuccess{\History}$ holds, a readout
$j\in\ResidualPositions$ with
$d_{\ResidualForest}(j)\le\DegreeCutoff$, and a candidate source
$i\in\ResidualPositions\setminus\{j\}$.  By
\eqref{app:eq:screen-success-event}, more than $\ColoringCount/2$
evaluations are $(i,j)$-separating.

Recall the binary dependence-test vote $\ScreenVote{\tau}{i}{j}$ in
\eqref{app:eq:screen-vote}: it is zero if
$\ColorHash{\tau}{i}=\ColorHash{\tau}{j}$, and otherwise
$\ScreenVote{\tau}{i}{j}
 =\mathbf 1\{\ScreenDiameter{\tau}{i}{j}>2\RowTVError\}$.
Here $\ScreenDiameter{\tau}{i}{j}$ is the observed TV diameter in
\eqref{app:eq:observed-screen-diameter}.
By \eqref{app:eq:provisional-row}, $i\in\CandidateRow{j}$ exactly when
$\sum_{\tau\in\ColoringIndexSet}\ScreenVote{\tau}{i}{j}
 >\ColoringCount/2$.

For a separating evaluation whose source and readout colors differ, define
the hidden exact rows
\begin{align}
 \ProbeExactRow{\tau}{\kappa}{i}{j}
 &:=\ExactRow{
   \ProbeState{\kappa}{\tau}{\ColorHash{\tau}{i}}
              {\ReadoutChunkOf{\tau}{j}}}{j},
 \qquad \kappa\in\ScreenColumnSet,
 \label{app:eq:hidden-screen-rows}\\
 \ExactRowFamily{\tau}{i}{j}
 &:=\{\ProbeExactRow{\tau}{\kappa}{i}{j}:
             \kappa\in\ScreenColumnSet\}.
 \label{app:eq:hidden-screen-row-family}
\end{align}

\emph{2. No false positives.}
Assume $\{i,j\}\notin E(\ResidualForest)$ and fix a separating $\tau$.
If $\ColorHash{\tau}{i}=\ColorHash{\tau}{j}$, then
$\ScreenVote{\tau}{i}{j}=0$ by \eqref{app:eq:screen-vote}.
Otherwise, the probe state in \eqref{app:eq:screen-probe-state} masks
$\ReadoutChunkOf{\tau}{j}$, so the other masked positions are exactly
$\ReadoutChunkOf{\tau}{j}\setminus\{j\}$.
Since $\ReadoutChunkOf{\tau}{j}
\subseteq\ColorClass{\tau}{\ColorHash{\tau}{j}}$, the first relation in
\eqref{app:eq:pairwise-local-isolation} gives
$\mathcal N_{\ResidualForest}(j)\cap\ReadoutChunkOf{\tau}{j}=\varnothing$.
The positions that can change across columns lie in the source class
$\ColorClass{\tau}{\ColorHash{\tau}{i}}$.
The second relation in \eqref{app:eq:pairwise-local-isolation} bounds its
intersection with $\mathcal N_{\ResidualForest}(j)$ by $\{i\}$;
this intersection is empty because $i\notin\mathcal N_{\ResidualForest}(j)$.
Thus all bank- and tail-column states agree on the committed history and
reveal every residual neighbor of $j$ at its fixed draft value.
Lemma~\ref{app:lem:residual-markov} gives
\begin{equation}
 \ProbeExactRow{\tau}{\TailColumn}{i}{j}
 =\ProbeExactRow{\tau}{\kappa}{i}{j}
 \quad\text{for every }\kappa\in\ScreenColumnSet.
 \label{app:eq:nonedge-exact-row-equality}
\end{equation}
Every observed row is within $\RowTVError$ of this common exact row by
\eqref{app:eq:oracle-tv-radius}; hence every pair of observed rows is within
$2\RowTVError$, and $\ScreenVote{\tau}{i}{j}=0$.
A strict majority of all evaluations are separating, so
$\sum_{\tau\in\ColoringIndexSet}\ScreenVote{\tau}{i}{j}<\ColoringCount/2$.
Thus $i\notin\CandidateRow{j}$ by \eqref{app:eq:provisional-row}, proving
\begin{equation}
 \CandidateRow{j}\subseteq\mathcal N_{\ResidualForest}(j).
 \label{app:eq:no-false-positive-row}
\end{equation}

\emph{3. No false negatives when the local vocabulary bank is nonempty.}
Now assume $\{i,j\}\in E(\ResidualForest)$ and $\BankSize{i}\ge1$.
For every separating $\tau$, the first relation in
\eqref{app:eq:pairwise-local-isolation} implies
$\ColorHash{\tau}{i}\ne\ColorHash{\tau}{j}$.
The second relation in \eqref{app:eq:pairwise-local-isolation} makes $i$
the only residual neighbor of $j$ in the probed source class.  Fix a complete
boundary $z\in\Vocab^{\PositionSet\setminus\{i,j\}}$ that agrees with the
history and assigns $\DraftToken{v}$ to every uncommitted coordinate
$v\ne i,j$.  Lemma~\ref{app:lem:fixed-draft-rows} gives
\begin{align}
 \ProbeExactRow{\tau}{\TailColumn}{i}{j}
 &=\BoundaryRow{i}{j}{\TailRepresentative{i}}{z},\notag\\
 \ProbeExactRow{\tau}{\kappa}{i}{j}
 &=\begin{cases}
 \BoundaryRow{i}{j}{\BankToken{i}{\kappa}}{z},
     &\kappa\le\BankSize{i},\\
 \BoundaryRow{i}{j}{\DraftToken{i}}{z},
     &\BankSize{i}<\kappa\le\MaxBankSize.
 \end{cases}
 \label{app:eq:isolated-boundary-row-family}
\end{align}
Every separating $\tau$ therefore contains the same bank-plus-tail exact
subfamily.  Its diameter is
\begin{equation}
 D^{\mathrm{ex}}_{\History}(i,j)
 :=\max_{a,b\in\TokenBank{i}\cup\{\TailRepresentative{i}\}}
 \dTV\bigl(\BoundaryRow{i}{j}{a}{z},
            \BoundaryRow{i}{j}{b}{z}\bigr).
 \label{app:eq:queried-exact-diameter}
\end{equation}
Choose columns $\kappa,\kappa'$ whose exact rows attain this finite
bank-plus-tail maximum.  Padding can add rows but cannot remove these
columns.  Applying \eqref{app:eq:oracle-tv-radius} to both rows gives
\begin{align*}
 \ScreenDiameter{\tau}{i}{j}
 &\ge\dTV\bigl(\ProbeObservedRow{\tau}{\kappa}{i}{j},
                  \ProbeObservedRow{\tau}{\kappa'}{i}{j}\bigr)\\
 &\ge\dTV\bigl(\ProbeExactRow{\tau}{\kappa}{i}{j},
                  \ProbeExactRow{\tau}{\kappa'}{i}{j}\bigr)\\
 &\quad-\dTV\bigl(\ProbeObservedRow{\tau}{\kappa}{i}{j},
                    \ProbeExactRow{\tau}{\kappa}{i}{j}\bigr)\\
 &\quad-\dTV\bigl(\ProbeObservedRow{\tau}{\kappa'}{i}{j},
                    \ProbeExactRow{\tau}{\kappa'}{i}{j}\bigr)\\
 &\ge D^{\mathrm{ex}}_{\History}(i,j)-2\RowTVError.
\end{align*}
If $D^{\mathrm{ex}}_{\History}(i,j)>4\RowTVError$, this is larger than
$2\RowTVError$, so $\ScreenVote{\tau}{i}{j}=1$ for every separating $\tau$.
The strict majority then puts $i$ in $\CandidateRow{j}$ by
\eqref{app:eq:provisional-row}.
Taking the contrapositive,
\begin{equation}
 i\notin\CandidateRow{j}
 \quad\Longrightarrow\quad
 D^{\mathrm{ex}}_{\History}(i,j)\le4\RowTVError.
 \label{app:eq:missed-queried-diameter}
\end{equation}
To rule out this missed edge, it suffices to show
$\DirectedResponse{i}{j}{z}\le\ScreenResolution<\EdgeSignal$,
contradicting (UEN\textcolor{red}{; \cref{app:eq:UEN}}).  We establish this full-response bound next.

\emph{Recovering the full response diameter.}
The tail representative $\TailRepresentative{i}$ and every token outside
the vocabulary bank $\TokenBank{i}$ lie in
$\TailSet{i}{\TailThreshold}$ by
Lemma~\ref{app:lem:preprocessing-certificates}.  Feasibility of
$\TailThreshold$ and (RT\textcolor{red}{; \cref{app:eq:RT}}) therefore give
\begin{equation}
 \max_{a,b\notin\TokenBank{i}}
 \dTV\bigl(\BoundaryRow{i}{j}{a}{z},
            \BoundaryRow{i}{j}{b}{z}\bigr)
 \le\min\{1,\ResponseConstant\TailThreshold^{\ResponseExponent}\}
 \le\TailTolerance.
 \label{app:eq:tail-diameter-at-screen-boundary}
\end{equation}
Suppose $i\notin\CandidateRow{j}$.  In the mixed case
$a\in\TokenBank{i}$ and $b\notin\TokenBank{i}$, the triangle inequality
through the tail representative gives
\begin{align*}
 &\dTV\bigl(\BoundaryRow{i}{j}{a}{z},\BoundaryRow{i}{j}{b}{z}\bigr)\\
 &\quad\le
 \dTV\bigl(\BoundaryRow{i}{j}{a}{z},
           \BoundaryRow{i}{j}{\TailRepresentative{i}}{z}\bigr)
 +\dTV\bigl(\BoundaryRow{i}{j}{\TailRepresentative{i}}{z},
             \BoundaryRow{i}{j}{b}{z}\bigr)\\
 &\quad\le 4\RowTVError+\TailTolerance.
\end{align*}
The first distance uses \eqref{app:eq:missed-queried-diameter}; the second
uses \eqref{app:eq:tail-diameter-at-screen-boundary}.  Interchanging $a,b$
handles the other mixed case.  Together with the bank--bank and tail--tail
bounds, this gives, for arbitrary $a,b\in\Vocab$,
\begin{equation}
 \dTV\bigl(\BoundaryRow{i}{j}{a}{z},
            \BoundaryRow{i}{j}{b}{z}\bigr)
 \le
 \begin{cases}
  4\RowTVError,&a,b\in\TokenBank{i},\\
  4\RowTVError+\TailTolerance,
    &\text{exactly one of $a,b$ lies in }\TokenBank{i},\\
  \TailTolerance,&a,b\notin\TokenBank{i}.
 \end{cases}
 \label{app:eq:screen-three-case-diameter}
\end{equation}
The maximum of the three upper bounds in
\eqref{app:eq:screen-three-case-diameter} is $\ScreenResolution$ by
\eqref{app:eq:diameter-screen-resolution}.  Hence a missed true edge would
have directed response at most $\ScreenResolution$, contradicting (UEN\textcolor{red}{; \cref{app:eq:UEN}}) and
$\ScreenResolution<\EdgeSignal$.

The same three comparisons also give the quantitative link used in the
intuition, without assuming that the edge was missed:
\begin{align*}
 \DirectedResponse{i}{j}{z}
 &\le\max\{D^{\mathrm{ex}}_{\History}(i,j),
           D^{\mathrm{ex}}_{\History}(i,j)+\TailTolerance,
           \TailTolerance\}\\
 &=D^{\mathrm{ex}}_{\History}(i,j)+\TailTolerance,\\
 \ScreenDiameter{\tau}{i}{j}
 &\ge D^{\mathrm{ex}}_{\History}(i,j)-2\RowTVError\\
 &\ge\DirectedResponse{i}{j}{z}-\TailTolerance-2\RowTVError\\
 &\ge\EdgeSignal-\TailTolerance-2\RowTVError>2\RowTVError.
\end{align*}
Here the three entries are the bank--bank, mixed, and tail--tail cases;
the last line uses (UEN\textcolor{red}{; \cref{app:eq:UEN}}) for this true edge.

\emph{4. Empty vocabulary banks and row selection.}
If $\BankSize{i}=0$, every source token lies in the certified tail, so
\eqref{app:eq:tail-diameter-at-screen-boundary} bounds the entire directed
response by
$\TailTolerance\le\ScreenResolution<\EdgeSignal$.  Such an $i$ cannot be
adjacent to $j$ by (UEN\textcolor{red}{; \cref{app:eq:UEN}}).  This also covers $\MaxBankSize=0$: then every local
vocabulary bank $\TokenBank{i}$ is empty, every returned row is defined to
be empty, and no residual edge
incident to a low-degree readout can exist.


\emph{5. Preserve the exact candidate row under public selection.}
We have proved that $\CandidateRow{j}$ equals the true neighbor set.  Its size
is at most $\DegreeCutoff$, so every admissible selector in
Definition~\ref{app:def:row-selector}, as applied in \AlgoBlockRef{aggregate},
preserves it by the nonoverflow identity condition.  Therefore
$\ScreenRow{j}=\mathcal N_{\ResidualForest}(j)$.
\end{proof}
